%=====================================================================
\chapter{Answers to Checkpoints}\label{app:answers}
%=====================================================================
\normalfigures

Answers to the three-question Checkpoint box in each chapter. Attempt them before
reading. Where an answer requires judgement, the reasoning matters more than the
conclusion.

\newcommand{\ansch}[1]{\subsection*{\color{primary}Chapter #1}}
\newenvironment{ans}{\begin{enumerate}[leftmargin=*, itemsep=4pt, topsep=3pt,
  label=\textbf{\arabic*.}]}{\end{enumerate}}

%---------------------------------------------------------------------
\ansch{1 --- The Data Science Landscape}
\begin{ans}
  \item It \emph{is} AI --- it performs a task that would otherwise need human
        clinical judgement. It is \emph{not} machine learning, because the rule was
        written by a person rather than learned from data. This is the distinction
        \dsref{ai-nesting} makes.
  \item (a) descriptive; (b) predictive; (c) diagnostic; (d) prescriptive.
  \item They are doing descriptive analytics. The distinction matters to whoever is
        paying because descriptive work is cheaper, faster and lower-risk than the
        predictive system the word ``AI'' implies --- and because expectations set at
        the wrong rung cause the project to be judged a failure for delivering
        exactly what was needed.
\end{ans}

\ansch{2 --- The Data Science Lifecycle}
\begin{ans}
  \item What the majority-class baseline scores. If 95\% of cases are negative,
        predicting ``negative'' always scores 95\%, so 94\% is \emph{worse} than
        doing nothing and 96\% is barely better. Ask for precision and recall.
  \item Because deployment changes the world the model was trained on: warned
        students behave differently, blocked attackers change tactics. The
        distribution shifts, so the cycle must restart (\chref{mlops}).
  \item Any feature recorded after or because of the readmission --- for example
        \texttt{days\_to\_readmission}, or a discharge-summary field completed on
        re-entry. It will not exist at prediction time, so the model cannot use it in
        production.
\end{ans}

\ansch{3 --- Data: Types, Sources and Quality}
\begin{ans}
  \item (a) nominal --- it is an identifier, not a quantity; (b) ordinal;
        (c) interval --- $0^\circ$C is not the absence of temperature;
        (d) ratio.
  \item The mean is dragged towards $-999$, and the variance is inflated
        enormously; both become meaningless. The mechanism is most likely
        \textbf{MNAR}, since a sensor fails to obtain a reading for reasons related
        to the condition being measured --- so encode the absence rather than
        imputing it.
  \item Aggregation. You lose the \emph{time within the week} at which attendance
        occurred, the identity of the specific sessions missed, and any
        within-week pattern. Whether that matters depends on the frame.
\end{ans}

\ansch{4 --- The Mathematical Toolkit}
\begin{ans}
  \item $\sqrt{3^2+4^2}=\sqrt{25}=5$.
  \item $\nabla L = 2w = -4$. The update is $w \leftarrow w - \alpha(-4)$, which
        \emph{increases} $w$ --- so it moves right, towards the minimum at 0.
  \item They are about the same subject (direction is nearly identical) but very
        different in magnitude --- typically a long document and a short one on the
        same topic. This is exactly the case where cosine similarity is right and
        Euclidean distance is misleading.
\end{ans}

\ansch{5 --- Describing Data}
\begin{ans}
  \item Strongly right-skewed. Publish the \textbf{median} of 12, with the IQR;
        the mean of 45 describes almost none of the observations.
  \item Mean $= 24/5 = 4.8$. Median $= 4$.
        $s^2 = \frac{(2-4.8)^2+(4-4.8)^2+(4-4.8)^2+(5-4.8)^2+(9-4.8)^2}{4}
             = \frac{7.84+0.64+0.64+0.04+17.64}{4} = 6.7$, so $s \approx 2.59$.
  \item Because the values are unordered category labels. ``Average protocol'' has
        no meaning, and a median requires an ordering that does not exist. Only the
        most frequent value is defined.
\end{ans}

\ansch{6 --- Probability and Uncertainty}
\begin{ans}
  \item About 2\%. Of 100{,}000 people, 100 have the condition and 95 test
        positive; 99{,}900 do not and about 4{,}995 test positive falsely. So
        $95/5090 \approx 1.9\%$. This is the base-rate fallacy
        (\dsref{base-rate}).
  \item No. Mutually exclusive means $P(A\cap B)=0$, but independence would require
        $P(A)P(B)=0.09$. Since $0 \neq 0.09$, they are strongly \emph{dependent}.
  \item No. It concerns the distribution of \emph{sample means}. Individual
        observations keep whatever distribution the population has, however large
        $n$ becomes.
\end{ans}

\ansch{7 --- Acquisition, Wrangling and Cleaning}
\begin{ans}
  \item The join key was not unique on the VLE side, so rows were duplicated ---
        probably multiple VLE accounts or enrolments per student. Check for
        duplicates on the key, decide the intended grain, and aggregate before
        joining.
  \item k-NN and k-means \textbf{require} scaling (both are distance-based).
        Random forest and gradient boosting do not --- they split one feature at a
        time, so units are irrelevant.
  \item (i) leakage from a feature computed after or during the failure window;
        (ii) duplicate rows spanning the train/test split, or splitting by row
        rather than by device; (iii) whether the test period is genuinely later in
        time than the training period.
\end{ans}

\ansch{8 --- Exploratory Data Analysis}
\begin{ans}
  \item Multicollinearity. Drop one of the pair (or use ridge). Keeping both adds
        almost no information and makes the coefficients unstable and
        uninterpretable.
  \item Yes --- a strong non-linear one. Pearson's $r$ measures only \emph{linear}
        association, which is why \dsref{scatter-patterns} exists.
  \item A heatmap with servers on one axis and days on the other, colour showing
        the failure rate. Forty line charts cannot be compared visually, and a
        single chart with forty lines is illegible; the heatmap makes both
        per-server and site-wide patterns visible at once.
\end{ans}

\ansch{9 --- Inferential Statistics}
\begin{ans}
  \item Not correct. The true mean is a fixed number, not a random one. The correct
        statement: if the sampling procedure were repeated many times, about 95\% of
        the intervals so constructed would contain the true mean.
  \item No. With $n=500{,}000$ almost any difference is statistically significant.
        A 0.03\% effect is almost certainly not worth acting on. Significance
        answers ``is it detectable?''; effect size answers ``is it worth anything?''
  \item $1 - 0.95^{15} \approx 54\%$. More than half the time you would find at
        least one ``significant'' result even if nothing were associated.
\end{ans}

\ansch{10 --- Relationships and Causality}
\begin{ans}
  \item The confounder is hot weather, which increases both ice cream sales and
        swimming. Explanation 3 in \dsref{four-explanations}.
  \item Because participation is not random: people who choose to participate differ
        systematically from those who do not (motivation, availability, need). A
        larger sample makes the biased estimate more precise, not more correct.
  \item Conditioning on a collider. Both traits increase the chance of being hired,
        so among the hired the two become negatively correlated. The association is
        an artefact of studying only successful applicants.
\end{ans}

\ansch{11 --- Machine Learning Fundamentals}
\begin{ans}
  \item Overfitting (high variance). Remedies: simplify the model or regularise;
        obtain more training data; use early stopping; reduce the feature count.
  \item Because every look at the test score followed by a change is a form of
        fitting to it. After twenty such iterations the test set has effectively
        become a validation set, and its score no longer estimates performance on
        genuinely unseen data.
  \item By \textbf{machine}. Readings from the same machine are highly correlated,
        so splitting by row would put near-identical rows in both train and test ---
        leakage that inflates the score.
\end{ans}

\ansch{12 --- Supervised Learning I: Regression}
\begin{ans}
  \item RMSE far exceeding MAE means the errors are dominated by a few very large
        ones --- the error distribution has a long tail. Investigate those cases
        specifically; they may be a distinct sub-population.
  \item No. The features are on different scales. If $x_1$ ranges 0--100 and $x_2$
        ranges 0--5, then $x_1$ contributes up to 50 and $x_2$ up to 60 --- roughly
        comparable. Standardise before comparing coefficients.
  \item Heteroscedasticity: the error variance grows with the prediction. Remedies:
        model $\log(y)$ instead, or use weighted least squares. The coefficients may
        still be usable, but the confidence intervals are wrong.
\end{ans}

\ansch{13 --- Supervised Learning II: Classification}
\begin{ans}
  \item $e^{1.1} \approx 3.0$. Each one-unit increase in the feature multiplies the
        odds of the positive class by about three, holding the other features
        constant.
  \item k-NN computes distances, so the feature with the largest numeric range
        dominates them. A tree splits on one feature at a time and only compares
        values within that feature, so its units are irrelevant.
  \item $G = 1 - (1^2 + 0^2) = 0$. The node is pure, so the tree stops splitting it
        and makes it a leaf.
\end{ans}

\ansch{14 --- Model Evaluation and Selection}
\begin{ans}
  \item No --- it is worse than useless. Predicting the majority class always would
        score 97\%. The model is 0.5 points \emph{below} doing nothing.
  \item Precision $= 60/500 = 12\%$. Recall $= 60/(60+20) = 60/80 = 75\%$.
  \item Ask for the precision--recall curve and for precision at a stated recall.
        With 0.2\% positives, ROC-AUC is dominated by the enormous negative class
        and can look excellent while the alerts are mostly wrong.
\end{ans}

\ansch{15 --- Feature Engineering and Ensembles}
\begin{ans}
  \item Bagging averages many independently-trained models, cancelling the
        uncorrelated part of their error --- it reduces \textbf{variance}. Boosting
        fits each new model to the previous one's residuals, progressively
        correcting systematic error --- it reduces \textbf{bias}.
  \item Otherwise every tree would choose the same dominant feature at the root and
        the trees would be nearly identical. Correlated models do not average away
        their errors, which defeats the purpose of bagging.
  \item Impurity-based importance is biased towards high-cardinality features, and
        an ID can perfectly separate training rows without generalising at all. It
        should not be a feature; use permutation importance to confirm.
\end{ans}

\ansch{16 --- Unsupervised Learning and Anomaly Detection}
\begin{ans}
  \item Because more centroids can only reduce the distance from each point to its
        nearest one; at $k=n$ inertia is zero. Minimising it therefore always
        chooses the largest $k$, so it must be combined with the elbow judgement,
        the silhouette score, and a domain reason.
  \item DBSCAN. k-means assumes roughly spherical clusters and forces every point
        into one, so the outliers you want identified would be absorbed into the
        nearest cluster and hidden.
  \item The unscaled feature with the largest variance --- typically the one with
        the biggest units --- dominated. PCA maximises variance, which is
        scale-dependent. Standardise first.
\end{ans}

\ansch{17 --- Neural Network Foundations}
\begin{ans}
  \item $z = (0.5)(2.0) + (-1.0)(1.0) + 0.3 = 1.0 - 1.0 + 0.3 = 0.3$.
        $\text{ReLU}(0.3) = 0.3$.
  \item Because the composition of linear functions is linear:
        $W_2(W_1x) = (W_2W_1)x$. A hundred linear layers compute what one linear
        layer computes.
  \item Overfitting. Stop training and restore the weights from the epoch with the
        lowest validation loss --- which is exactly what early stopping automates.
\end{ans}

\ansch{18 --- Deep Architectures}
\begin{ans}
  \item The kernel has $3\times3 = 9$ weights (plus one bias). A fully connected
        layer mapping $32\times32 = 1024$ inputs to a similar number of outputs would
        need on the order of a million weights. This is the parameter saving that
        makes CNNs trainable.
  \item Attention connects any two positions in a single step, so the gradient path
        between them has length 1. An LSTM must propagate information through 99
        intermediate steps, and the gradient degrades along the way.
  \item A model trained from scratch on 600 images will memorise them. A pre-trained
        network already detects edges, textures and shapes --- most of what any
        visual task needs --- so only the final layers must be fitted, which 600
        images can support. It also trains in minutes rather than days.
\end{ans}

\ansch{19 --- Sequences, Time Series and Streaming}
\begin{ans}
  \item A centred window includes values \emph{after} time $t$, which will not exist
        when the model runs. The feature encodes the future, so validation scores are
        inflated and production performance collapses.
  \item Not necessarily. Decompose first: if traffic drops every 25 December, the
        seasonal component accounts for it and the residual is unremarkable. Only a
        drop beyond the expected seasonal pattern is an anomaly.
  \item (i) Did an upstream data source change or break --- check input distributions
        and row counts. (ii) Did the world change --- a new system, a policy change,
        a new attacker technique. A sudden cliff is almost always an upstream change
        rather than gradual model decay.
\end{ans}

\ansch{20 --- Text and NLP}
\begin{ans}
  \item $\log(N/N) = \log 1 = 0$, so its TF-IDF score is zero regardless of how often
        it appears. IDF removes uninformative words automatically.
  \item ``The system is \emph{not} vulnerable'' becomes ``system vulnerable'' once
        \emph{not} is removed as a stop word --- an exact reversal. Bigrams or a
        curated stop-word list avoid this.
  \item Because bag-of-words dimensions correspond to exact strings, and ``server''
        and ``host'' are simply different strings. Dense embeddings, which place
        words used in similar contexts near each other, fix it.
\end{ans}

\ansch{21 --- Generative AI and LLMs}
\begin{ans}
  \item The model samples the most \emph{plausible} next token given the context; it
        has no separate fact store to check against and no representation of its own
        uncertainty about the world. A fabricated citation is produced by the same
        process, with the same fluency, as a correct one.
  \item RAG. Fine-tuning teaches behaviour and format, not facts, and a policy that
        changes would require retraining every time. Retrieval also supplies
        citations, so the answer can be checked.
  \item Temperature 0. The model then always takes the highest-probability token,
        giving deterministic output --- the right setting for extraction,
        classification and anything that must be consistent between runs.
\end{ans}

\ansch{22 --- Agentic AI}
\begin{ans}
  \item \emph{Agent:} investigating an unfamiliar alert, where the next step depends
        on what the previous one revealed and several tools may be needed.
        \emph{Script:} generating the weekly report from a fixed query --- known,
        repeated, and better served by something testable and predictable.
  \item Because retrieved content is attacker-controllable text, and an LLM has no
        reliable boundary between data and instructions. A classifier only outputs a
        label; an agent with tools can be induced to \emph{act} on injected
        instructions.
  \item A maximum-iteration cap and a repeated-action detector. A token or spend
        budget would also have stopped it.
\end{ans}

\ansch{23 --- Data Engineering at Scale}
\begin{ans}
  \item Idempotent means re-running a task produces the same result rather than
        duplicating or compounding its effect. It matters because failures are
        normal: a job that half-completed must be safe to re-run without producing
        double-counted rows.
  \item No. 300 MB fits comfortably in memory; pandas will be faster to write, faster
        to run and far easier to debug. Distributed tooling on small data adds days
        of complexity for negative benefit.
  \item \textbf{Completeness}, caught by the \textbf{row-count / volume check}
        against a rolling median. Every individual row is valid, so type, range and
        null checks all pass --- only the count reveals it.
\end{ans}

\ansch{24 --- Cloud and Edge Computing}
\begin{ans}
  \item \textbf{You are}, entirely. Under shared responsibility the provider secures
        the infrastructure; access configuration and data are the customer's side of
        the line. Nearly all publicised ``cloud breaches'' are of this kind.
  \item (i) It pins the system libraries, interpreter version and OS packages, not
        just Python packages; (ii) it is a single immutable artefact that can be
        re-run identically months later, whereas \texttt{pip install -r} resolves
        differently over time.
  \item The \textbf{edge}. Cloud round-trip latency alone is 50--500 ms, so the
        budget cannot be met; and a network outage would remove all protection.
        Safety-critical latency is not negotiable by improving the model.
\end{ans}

\ansch{25 --- Data Science for IoT}
\begin{ans}
  \item The sensor is stuck --- a common failure mode that produces a plausible
        constant. The check is \textbf{rolling variance too low}: a real process
        always shows some noise.
  \item Because the decision is made per device per maintenance cycle, not per
        reading. At raw grain the class balance is absurd, nearly every row is
        healthy, and the effective sample size is the number of failures, not the
        number of rows.
  \item No. A 24-hour warning cannot trigger a three-week procurement, so the
        prediction cannot change the outcome. A model with lower recall but a
        four-week lead time would be more valuable.
\end{ans}

\ansch{26 --- Data Science for Cybersecurity}
\begin{ans}
  \item False positives $\approx 500{,}000 \times 0.01 = 5{,}000$; true positives
        $\approx 10$. About \textbf{5{,}010 alerts/day}, precision $\approx 0.2\%$.
        With four analysts (roughly 160 alerts/day) it is 30$\times$ over capacity
        and \textbf{not deployable}.
  \item Because packets per second is directly and cheaply controllable by the
        attacker --- they simply slow down. The fraction of connections to closed
        ports is intrinsic to scanning: avoiding it requires already knowing which
        ports are open, which is what the scan is for.
  \item \textbf{Poisoning.} Controls: screen the training set with an anomaly
        detector before use; require human approval to promote a retrained model;
        hold out a trusted benchmark set that a poisoned model would fail.
\end{ans}

\ansch{27 --- MLOps}
\begin{ans}
  \item Training and serving compute features with different code, so they disagree
        subtly --- for example a different timezone or null convention. Example:
        training uses \texttt{fillna(median)} and serving uses \texttt{fillna(0)}.
        A \textbf{feature store} eliminates the class of bug by computing features
        once for both paths.
  \item \textbf{Layers 1 and 2} --- input distributions and prediction
        distributions. Layer 1 detects broken pipelines, changed upstream schemas and
        data drift; layer 2 detects a model whose behaviour has shifted, e.g. flagging
        12\% where it used to flag 3\%.
  \item Because offline metrics do not capture training--serving skew, latency,
        or the effect on real behaviour. Shadow and canary deployment test those on
        live traffic before the new model drives any decision.
\end{ans}

\ansch{28 --- Responsible AI and Data Ethics}
\begin{ans}
  \item No. The model will reconstruct gender from proxies --- programme, device
        type, activity timing --- and you have now lost the ability to \emph{measure}
        the disparity. Test by trying to predict the protected attribute from the
        remaining features; an AUC well above 0.5 confirms proxies exist.
  \item No. This is the impossibility result: with different base rates, a calibrated
        model necessarily flags the higher-base-rate group more often. Forcing equal
        flag rates breaks calibration.
  \item A \textbf{counterfactual} explanation: ``you were flagged mainly because you
        have not submitted since week 2; submitting this week's work would remove the
        flag.'' It is the only form that tells the person what they can do.
\end{ans}

\ansch{29 --- Research Methods}
\begin{ans}
  \item For example: ``Among mid-sized food distributors in Pakistan, does adopting a
        blockchain-based provenance system reduce the mean time to trace a
        contaminated batch, compared with the existing paper-based process?''
        --- population, intervention, outcome, comparison and scope all stated.
  \item \textbf{Construct validity}. Clicks measure activity, not learning; a student
        may click without engaging or engage without clicking. Strengthen it with
        multiple indicators --- clicks, submissions, time on task and self-report ---
        and report their agreement.
  \item \emph{Reproducible}: the same data and code give the same result.
        \emph{Replicable}: new data and the same method give the same conclusion.
        The first is an engineering property; the second is what science requires.
\end{ans}

\ansch{30 --- Managing Data Science Projects}
\begin{ans}
  \item Feasibility is unknown at the start. In software you know the feature can be
        built and estimate how long; here the signal may simply not be in the data,
        and no planning discipline can determine that in advance.
  \item Commit to a \emph{gate}: a date, a deliverable and a decision --- ``in three
        weeks we will report whether the data supports 70\% recall at a 20\% flag
        rate.'' Explain that a date is within your control and a result is not, and
        that a fabricated accuracy figure would be worse than useless for planning.
  \item No --- it is the process working. Three weeks of effort established that the
        project was not viable, instead of nine months. Naming the gates in advance
        is what makes this an outcome rather than a defeat.
\end{ans}

\ansch{31 --- Communicating Data Science}
\begin{ans}
  \item ``Of every 100 cases we flag, about 82 are genuine.''
  \item \emph{Which decision does this chart change, and who makes it?} Anything that
        cannot change an action comes out.
  \item Because an unstated limitation becomes an assumed capability. The reader will
        rely on the system in situations you know it cannot handle, and will discover
        the gap in production --- with less context, less goodwill, and possibly
        someone harmed.
\end{ans}

\ansch{32 --- Capstone}
\begin{ans}
  \item (i) Is it leakage? Check every feature against the prediction-time cutoff and
        check the split for entity or time overlap. (ii) What is the majority-class
        baseline? At 6\% positives it is 94\%, so 0.97 is a much smaller gain than it
        appears. Only then look at precision and recall.
  \item (i) Labels appear but rows collapse from millions to thousands, so
        unsupervised anomaly detection gives way to supervised learning with careful
        validation. (ii) The entity becomes a person, so interpretability and a
        contest route become requirements rather than preferences. (iii) Ground-truth
        lag stretches to a full term, so input and prediction monitoring must carry
        the alerting load.
  \item The \textbf{one-page decision brief}. Its first two sentences must state what
        the system can do, in the sponsor's units, and what you recommend doing about
        it --- not the method.
\end{ans}
